% part: intuitionistic-logic
% chapter: introduction
% section: constructive-reasoning

\documentclass[../../../include/open-logic-chapter]{subfiles}

\begin{document}

\olfileid{int}{int}{cr}

\section{নির্মাণমূলক যুক্তিবিচার}

মোডাল অপারেটর বা দ্বিতীয়-ক্রমের পরিমাণসূচক যোগ করে ধ্রুপদি যুক্তিবিদ্যাকে
সম্প্রসারিত করার বিপরীতে, স্বজ্ঞাবাদী যুক্তিবিদ্যা ধ্রুপদি যুক্তিবিদ্যাকে
সীমিত করে বলে এটি ``অধ্রুপদি''। ধ্রুপদি যুক্তিবিদ্যা নানা দিক থেকে
\emph{অনির্মাণমূলক}। স্বজ্ঞাবাদী যুক্তিবিদ্যার উদ্দেশ্য হলো এক ধরনের
নির্মাণমূলক গণিতের বৈশিষ্ট্যপূর্ণ, আরও ``নির্মাণমূলক'' যুক্তিবিচারকে
ধরা। নিচের উদাহরণগুলি এর অন্তর্নিহিত কয়েকটি প্রেরণা বোঝাতে পারে।

ধরো, কেউ দাবি করল যে সে এমন একটি স্বাভাবিক সংখ্যা~$n$ নির্ণয় করেছে
যার ধর্ম হলো: $n$ জোড় হলে রিমান অনুমান সত্য, আর $n$ বিজোড় হলে
রিমান অনুমান মিথ্যা। দারুণ খবর!{} রিমান অনুমান সত্য কি না, তা গণিতের
বড় অমীমাংসিত প্রশ্নগুলির একটি; আর মনে হচ্ছে সমস্যাটি একটি গণনার
প্রশ্নে নামিয়ে আনা হয়েছে---অর্থাৎ, একটি নির্দিষ্ট সংখ্যা জোড় কি
না তা নির্ণয় করলেই হবে।

$n$-এর সেই জাদুকরি মানটি কী? লোকটি এভাবে বর্ণনা করল: রিমান অনুমান
সত্য হলে $n$ হলো $2$, আর তা না হলে $3$।

রেগে গিয়ে তুমি টাকা ফেরত চাইলে। ধ্রুপদি দৃষ্টিভঙ্গি থেকে উপরের
বর্ণনাটি সত্যিই $n$-এর একটি অনন্য মান নির্ধারণ করে; কিন্তু তুমি আসলে
চাও $n$-এর মানটি \emph{স্পষ্টভাবে} দেওয়া হোক।

আরেকটি, সম্ভবত কম কৃত্রিম উদাহরণ নেওয়া যাক। আমরা জানি, একটি অমূলদ
সংখ্যাকে মূলদ ঘাতে তুলে মূলদ ফল পাওয়া সম্ভব। যেমন,
$\sqrt{2}^2 = 2$। যেটি কম পরিষ্কার তা হলো, কোনো অমূলদ সংখ্যাকে
\emph{অমূলদ} ঘাতে তুলেও মূলদ ফল পাওয়া সম্ভব কি না। নিচের উপপাদ্যটি
ইতিবাচক উত্তর দেয়:

\begin{thm}
এমন অমূলদ সংখ্যা $a$ ও $b$ আছে যেন $a^b$ মূলদ।
\end{thm}

\begin{proof}
$\sqrt{2}^{\sqrt{2}}$ বিবেচনা করো। এটি মূলদ হলে কাজ শেষ: $a = b =
\sqrt{2}$ নেওয়া যায়। তা না হলে এটি অমূলদ। তখন
\[
(\sqrt{2}^{\sqrt{2}})^{\sqrt{2}} = \sqrt{2}^{\sqrt{2} \cdot
  \sqrt{2}} = \sqrt{2}^2 = 2,
\]
যা মূলদ। তাই এই ক্ষেত্রে $a$ হিসেবে
$\sqrt{2}^{\sqrt{2}}$ এবং $b$ হিসেবে~$\sqrt 2$ নাও।
\end{proof}

এটি কি বৈধ প্রমাণ? অধিকাংশ গণিতবিদ মনে করেন, হ্যাঁ। কিন্তু এখানেও
একটু অসন্তোষ থেকে যায়: একটি নির্দিষ্ট ধর্মবিশিষ্ট বাস্তব সংখ্যার
একটি জোড়ার অস্তিত্ব আমরা প্রমাণ করেছি, অথচ সেটি \emph{কোন} জোড়া তা
বলতে পারিনি। একই ফল এমনভাবে প্রমাণ করা সম্ভব যাতে প্রমাণের মধ্যেই
$a$, $b$ জোড়াটি \emph{দেওয়া} থাকে: $a = \sqrt{3}$ এবং $b = \log_3
4$ নাও। তাহলে
\[
a^b = \sqrt{3}^{\log_3 4} = 3^{1/2 \cdot \log_3 4} = (3^{\log_3
  4})^{1/2} = 4^{1/2}= 2,
\]
কারণ $3^{\log_3 x} = x$।

প্রথম প্রমাণটির মতো পদক্ষেপ নিষিদ্ধ এমন এক ধরনের যুক্তিবিচার ধরার
জন্য স্বজ্ঞাবাদী যুক্তিবিদ্যা নির্মিত। $!A(x)$ পরিতৃপ্তকারী কোনো
$x$-এর অস্তিত্ব প্রমাণ করার অর্থ হলো, দ্বিতীয় প্রমাণের মতো একটি
নির্দিষ্ট~$x$ এবং সেটি $!A$ পরিতৃপ্ত করে তার প্রমাণ দিতে হবে। $!A$
অথবা $!B$ সত্য প্রমাণ করার জন্য দুটির কোনো একটিকে প্রমাণ করতে পারতে
হবে।

আনুষ্ঠানিকভাবে বলতে গেলে, ধ্রুপদি যুক্তিবিদ্যার !!a{derivation}
ব্যবস্থাকে একটি নির্দিষ্ট উপায়ে সীমিত করলে স্বজ্ঞাবাদী যুক্তিবিদ্যা
পাওয়া যায়। গাণিতিক দৃষ্টিতে এগুলি কেবল আনুষ্ঠানিক নিগমনব্যবস্থা;
কিন্তু, আগেই বলা হয়েছে, এগুলির উদ্দেশ্য এক ধরনের গাণিতিক যুক্তিবিচার
ধরা। কেউ একে গণিত সম্পর্কে একটি বিশেষ দার্শনিক দৃষ্টিভঙ্গিতে (যেমন,
Brouwer-এর স্বজ্ঞাবাদে) সমর্থিত যুক্তিবিচার বলে নিতে পারেন; কেউ একে
আরও ``সুনির্দিষ্ট'' ও সন্তোষজনক গাণিতিক যুক্তিবিচার বলে নিতে পারেন
(Bishop-এর নির্মাণবাদ অনুসারে); আবার আনুষ্ঠানিক বিবরণটি অনানুষ্ঠানিক
প্রেরণাকে সত্যিই ধরে কি না, তা নিয়েও বিতর্ক করা যায়। কিন্তু আমাদের
দার্শনিক অবস্থান যা-ই হোক, স্বজ্ঞাবাদী যুক্তিবিদ্যাকে আনুষ্ঠানিকভাবে
উপস্থাপিত যুক্তিবিদ্যা হিসেবে অধ্যয়ন করা যায়; এবং নানা কারণেই অনেক
গাণিতিক যুক্তিবিদ তা করতে আগ্রহী।

\end{document}
